\chapter{Central Timor}\label{ch:CT}
The \tsc{\ili{Central} Timor} subgroup contains several languages of
central Timor: \ili{Welaun} [wlh],
\ili{Kemak} [kem], \ili{Tokodede} [tkd],
and the \tsc{Mambae} language/dialect cluster [mgm].
The locations of these languages are shown in \frf{05-01}.
Evidence for \tsc{\ili{Central} Timor} comes primarily from shared *ŋ > **g.
Supporting evidence comes from *R > {\0},
**r; *ma-p > **mp;
and *ə > \it{a} in final syllables.
These changes are discussed in more detail in \srf{05-02} below.

\begin{figure}\centering
	\includegraphics[width=\columnwidth]{CT.pdf}
	\fcp{\tsc{\ili{Central} Timor} languages}{05-01}
\end{figure}

All of these languages have had contact with \ili{Tetun} [tet] which
is socially dominant in central
and eastern Timor and has functioned as a regional lingua franca for centuries.
This function is now institutional in Timor-Leste,
as \ili{Tetun} Dili [tdt] -- a \ili{Tetun}-based creole -- is one of the official languages.

Among all \tsc{\ili{Central} Timor} languages,
\ili{Welaun} shows particularly strong signs of influence from \ili{Tetun}
[tet], specifically West \ili{Tetun} (\ili{Foho} \ili{Tetun}
and \ili{Fehan} \ili{Tetun}), which is spoken in the central Timor region
straddling the modern political border.
\ili{Welaun} shares many sound changes,
structural features, and lexical items with \ili{Tetun} to such an
extent that it is often impossible to determine whether a particular
lexical item is a native \ili{Welaun} term or a \ili{Tetun} loan.
\citet[35f]{Edwards2019} found that \ili{Welaun} showed more lexical
similarity to \ili{Tetun} than to languages of the \tsc{\ili{Central} Timor}
subgroup, with similarity between \ili{Welaun}
and \ili{Tetun} ranging from 55--58{\%}.
Indeed, nearly a third of \ili{Welaun} vocabulary on a basic word list
is identical to \ili{Tetun} (ignoring productive morphology).

In cases where a \ili{Welaun} word is formally
and semantically identical to \ili{Tetun},
but does not follow the regular \ili{Welaun} sound changes,
it is straightforward to identify it as a loan.
Thus, for instance, \ili{Welaun} *duha > \it{rua} ‘two’ does not show
regular *d > \it{l} (see \srf{05-02})
and is thus probably a \ili{Tetun} loan.
Similarly, \ili{Welaun} \it{nani} ‘swim’ does not show regular *ŋ >
**g > \it{k} (compare \ili{Kemak}
and \ili{Tokodede} \it{nagi} ‘swim’)
and is thus probably a \ili{Tetun} loan.

However, in cases where \ili{Welaun}
and \ili{Tetun} share the same sound changes, it is often impossible
to distinguish between inheritances and loans.
For instance, both \ili{Welaun}
and \ili{Tetun} have undergone initial *p > \it{h}.
There are at l\ili{east} two \ili{Welaun} words which are different to their
\ili{Tetun} equivalents and establish this as a sound change in \ili{Welaun}: *əpat > **pəat
> \it{hoat} ‘four’ and *pənuq > \it{honu} ‘full’
(\ili{Tetun} \it{haat} ‘four’ and \it{konu} ‘full’).
However, other words such as *pitu > \it{hitu} ‘seven’
and *puqun > \it{huun} ‘tree trunk,
source, base’ are identical in \ili{Welaun} and \ili{Tetun}.

\begin{table}\tcp{\tsc{\ili{Central} Timor} *p = **p}{05-01}
%\begin{adjustbox}{width=\textwidth}
	\begin{threeparttable}\st{2pt}
	\begin{tabular}{ll@{ }llllll}\lsptoprule
gloss	&	{\hr}four	&	{\hr}how m.?&	{\hr}trunk	&	{\hr}navel, centre	&	{\hr}fire	&	{\hr}teeth\\
PMP	&	\hr*ə\tbr{p}at	&	\hr*\tbr{p}ija	&	\hr*\tbr{p}uqun	&	\hr*\tbr{p}usəj	&	\hr*ha\tbr{p}uy	&	\hr*ni\tbr{p}ən\\
P.-C. Timor	&	**\tbr{p}əat	&	**\tbr{p}ija	&	**\tbr{p}uun	&	**\tbr{p}usar	&	**a\tbr{p}i	&	**ni\tbr{p}an\\ \midrule
\ili{Welaun}	&	{\hr\hr}\tbr{h}oat	&	{\hr\hr}\tbr{h}ira	&	{\hr\hr}\tbr{h}uun	&	{\hr\hr}\tbr{h}usar-aat	&	{\hr\hr}haʔi	&	\hr\cg{(kiis-aat)}\\
\ili{Kemak}	&	{\hr\hr}\tbr{p}aat	&	{\hr\hr}\tbr{p}ila	&	{\hr\hr}ai \tbr{p}uun	&	{\hr\hr}\tbr{p}usr-aar\su{Δ}	&	{\hr\hr}a\tbr{p}i	&	{\hr\hr}ni\tbr{p}a-r\\
\ili{Tokodede}	&	{\hr\hr}\tbr{p}aat	&	{\hr\hr}\tbr{p}iil	&	{\hr\hr}\tbr{p}uu	&	{\hr\hr}ku\tbr{p}usa	&	{\hr\hr}a\tbr{p}i	&	{\hr\hr}ni\tbr{p}a\\
NW. \ili{Mambae}	&	{\hr\hr}\tbr{p}aat	&	{\hr\hr}ar\tbr{p}iil\su{‡}		&	{\hr\hr}\tbr{p}uu	&		&	{\hr\hr}ae\tbr{p}-a	&	{\hr\hr}ni\tbr{p}an\\
C. \ili{Mambae}	&	{\hr\hr}\tbr{f}aat	&	{\hr\hr}ar\tbr{f}iil\su{‡}		&	{\hr\hr}\tbr{f}uu	&		&	{\hr\hr}ai\tbr{f}-a	&	{\hr\hr}ni\tbr{f}an\\
S. \ili{Mambae} (H.)\su{†}	&	{\hr\hr}\tbr{p}aat	&	{\hr\hr}ar\tbr{p}iil\su{‡}		&		&	{\hr\hr}\tbr{p}usa	&	{\hr\hr}aa\tbr{p}	&	{\hr\hr}ni\tbr{p}a\\
S. \ili{Mambae} (B.)	&	{\hr\hr}\tbr{f}aat	&	{\hr\hr}ar\tbr{f}iil\su{‡}		&	{\hr\hr}\tbr{f}uu	&	{\hr\hr}\tbr{f}usu	&	{\hr\hr}ae\tbr{f}	&	{\hr\hr}ni\tbr{f}a\\
		\lspbottomrule
	\end{tabular}
		\begin{tablenotes}\tnsize
			\item[†]	{South \ili{Mambae} (H.) data is from the Hatu-Udo subdistrict.
								South \ili{Mambae} (B.) data is from Betano village in the Same subdistrict.}
			\item[‡]	{\tsc{Mambae} \it{arpiil, arfiil} = ‘when?’}
			\item[Δ]	{\ili{Kemak} \it{pusr-aar} is from \ili{Atabae} \ili{Kemak} \citep{Elliot2020}.}
		\end{tablenotes}
	\end{threeparttable}
%\end{adjustbox}
\end{table}

There is no known strong phonological evidence for placing \tsc{\ili{Central}
Timor} with other languages of the region below the level of MP.
Evidence that \tsc{\ili{Central} Timor} is not a member of \tsc{Timor-Babar} -- its
closest Austronesian geographic neighbour -- comes from
the fact that *p > \it{h,} {\0} cannot be reconstructed to Proto-\ili{Central} Timor.
Instead, *p is retained unchanged in most words in \ili{Tokodede},
\ili{Kemak}, Northwest \ili{Mambae}, and some varieties of South \ili{Mambae}.
\ili{Central} \ili{Mambae} and other varieties of South \ili{Mambae} have *p > \it{f}.
\ili{Welaun} has initial *p > \it{h}.
(The reflex of word-medial *p in \ili{Welaun} is unclear due to lack
of cognates which are not identical to \ili{Tetun}.\footnote{
		There is only one putative reflex of PMP *p intervocalically which
		is not identical to West \ili{Tetun}: *tapi > \it{tabi} ‘winnow’
		(West \ili{Tetun} \it{tataʔek} ‘winnow’).
		This would show *p > \it{b} intervocalically.
		Other \ili{Welaun} words with reflexes of intervocalic *p are identical
		to West \ili{Tetun} cognates, from which they might be borrowed.
		West \ili{Tetun} has *p > **h > \it{ʔ} word medially.})
Examples are given in \trf{05-01}.

There are also a small number of words which show *p > \it{h}
in all \tsc{\ili{Central} Timor} languages.
One of these is \ili{Welaun},
\ili{Kemak}, and \ili{Central} \ili{Mambae} *pitu > \it{hitu} ‘seven’.
Another example is \it{mehi} ‘dream’
(ultimately from PMP *ma-hipi)
in several \ili{Kemak} varieties and South \ili{Mambae}.
That this is a borrowing from East \ili{Tetun} \it{mehi} (West \ili{Tetun}
\it{meʔi}) is indicated by \ili{Kutubaba} \ili{Kemak} \it{mə̆niip}
and \ili{Leosibe} \ili{Kemak} \it{mə̆nipi} from PMP *ma-h\<in\>ipi ‘dream’
in which medial *p is retained unchanged.\footnote{
		Throughout this chapter reduced ante-penultimate vowels in \ili{Kemak} are transcribed ‹ə̆›.
		Most such vowels are reflexes of reconstructed vowels,
		but the synchronic analysis of these vowels is unclear.
		They could be epenthetic insertions to break up consonant clusters,
		or they could be underlying vowels which are reduced due to being
		in pre-stressed ante-penultimate position.}

\scp{\tsc{Central Timor} internal diversity}{05-01}

Each of the \tsc{Central Timor} “languages” are internally diverse
with multiple dialects and/or varieties.
Kemak speakers report about at least a dozen different named
dialects of their language including: Atsabe,
Diirbatin, Kailaku, Kutubaba, Lemia,
Leolima, Leosibe, Marobo, Oho, and Saneri.
Data for about ten of these varieties is available,
with lexical similarity between any two varieties ranging from
about 76{\%} to 88{\%}.\footnote{
		Kutubaba Kemak is
		the most divergent known variety of Kemak.
		Kutubaba Kemak is spoken in Kenebibi village on the north coast
		of the easternmost Kemak-speaking area in \frf{05-01}.}

\citet[i]{LSG2006} lists three different varieties of Tokodede:
Bazartete Tokodede, Liquiçá Tokodede, and Maubara Tokodede.
Lexical differences are found between all three varieties as
well as phonological differences,
the most significant of which is Maubara Tokodede /v/ corresponding
to other Tokodede /h/.
One example is *babuy > Maubara Tokodede \it{vavi}, other Tokodede \it{hahi} ‘pig’.

\citet{Fogaca2017} shows that varieties of \tsc{Mambae} fall
into three groupings on the basis of lexical similarity: Northwest
Mambae, Central Mambae, and South Mambae.
The South Mambae grouping is most divergent with lexical similarity
compared to other varieties of \tsc{Mambae} ranging from 63{\%}--70{\%}.\footnote{
		The degree of lexical similarity between Northwest Mambae
		and South Mambae is similar to the degree of lexical similarity
		between Northwest Mambae and Tokodede, which is about 61{\%} \citep[35]{Edwards2019}.}
\citeauthor{Fogaca2017}’s figures summarising the degree of lexical similarity
between different varieties of \tsc{Mambae} are given in \trf{05-02}.


\begin{table}
\caption[\tsc{Mambae} lexical similarity (Fogaça 2017: 82)]
{\tsc{Mambae} lexical similarity \citep[82]{Fogaca2017}\su{†}}\label{tab:05-02}
%\begin{adjustbox}{width=\textwidth}
	\begin{threeparttable} %\st{2pt}
	\begin{tabular}{llllllllllll}
\hhline{---------}
\mc{3}{|l}{Bazartete}	&		&		&		&		&		&	\mc{1}{c|}{}	&&				 						\\
\mc{1}{|l}{81}	&	\mc{3}{l}{Hatulia}				&	\mc{5}{r|}{\tsc{Northwest Mambae}}	&	\\
\mc{1}{|l}{78}	&	81	&	\mc{2}{l}{Railaco}	&		&		&		&		&	\mc{1}{c|}{}	&&	\\ \hhline{-----------}
74	&	75	&	76	&	\mc{3}{|l}{Laulara}	&		&	& & &	\mc{1}{c|}{} \\%&		&		&		&	\\
74	&	76	&	76	&	\mc{1}{|l}{85}	&	\mc{3}{l}{Aileu Vila}	&		&		\mc{3}{r|}{\multirow{2}{*}{\tsc{Central Mambae}}}	\\
71	&	73	&	70	&	\mc{1}{|l}{81}	&	83	&	\mc{5}{l}{Hatu-Builico}	&			\mc{1}{c|}{}\\
69	&	69	&	70	&	\mc{1}{|l}{82}	&	80	&	82	&	\mc{4}{l}{Liquidoe}	&		\mc{1}{c|}{} \\ \hhline{~~~---------}
63	&	65	&	65	&	68	&	70	&	68	&	66	&	\mc{3}{|l}{Hatu-Udo}	&		&	\mc{1}{c|}{}\\
63	&	64	&	64	&	67	&	68	&	68	&	67	&	\mc{1}{|l}{82}	&	\mc{2}{l}{Letefoho}	&	\mc{2}{r|}{\tsc{South Mambae}}\\
63	&	66	&	63	&	66	&	68	&	67	&	65	&	\mc{1}{|l}{82}	&	80	&	Betano	& & \mc{1}{c|}{}\\ \hhline{~~~~~~~-----}
		\lspbottomrule
	\end{tabular}
		\begin{tablenotes}\tnsize
			\item[†] {Names of different varieties are from subdistricts where they are spoken.
								Letefoho and Betano are \it{sucos} (“villages”) in the subdistrict of Same.}
			%\item[‡]
			%\item[Δ]
			%\item[‖]
			%\item[¶]
		\end{tablenotes}
	\end{threeparttable}
%\end{adjustbox}
\end{table}

Because this chapter is focused on the internal relationships
of the languages of \tsc{Central Timor},
it necessarily highlights the similarities between these languages.
To give a sense of the structural,
phonological, and lexical diversity which is encompassed under
the label “\tsc{Mambae}” numerals in five different varieties
of \tsc{Mambae} are given in \trf{05-03}.

Structural diversity can be seen in the way the numerals 6--9
and multiples of 10 are formed.
While both Northwest Mambae
and South Mambae have complex forms for the numerals 6--9 (see
\srf{02-03-01-04}), they use different constructions.
Northwest Mambae uses \it{hoho} (no known independent meaning
--- also used in Tokodede for numerals 6--9).\footnote{
		Maubara Tokodede uses \it{vovo} for numerals 6--9 pointing to earlier **bobo.}
while South Mambae uses \it{liim nai} ‘five \it{nai}’.
Central Mambae differs again with mono-morphemic forms which
are ultimately from PMP.
For the multiples of ten, three different constructions are found:
Northwest Mambae \it{guul} X,
Central Mambae X \it{nulo} (or metathesised \it{nuul}),
and South Mambae \it{saguul haet} X ‘ten times X’ or just \it{haet} ‘times’.
The numerals in \trf{05-03} also give a sample of the phonological
diversity within the \tsc{Mambae} cluster.
Northwest Mambae /p/ corresponds to /f/ in Central Mambae
and either /p/ or /f/ in different varieties of South Mambae.
Similarly, Central Mambae has /k/ where other varieties have /g/.

\begin{table}
	\centering\tcp{Numerals in five varieties of \tsc{Mambae}}{05-03}
	\st{3.9pt}\begin{tabular}{llllll}\lsptoprule
	&	NW. Mambae	&	C. Mambae	&	C. Mambae	&	S. Mambae	&	S. Mambae\\
	&	Railaco	&	Liquidoe	&	Hatu-Builico	&	Hatu-Udo	&	Betano\\ \midrule
1	&	iid	&	iid	&	iid	&	iid	&	iid\\
2	&	ruu	&	ruu	&	ruu	&	rua	&	ruu\\
3	&	teul	&	teul	&	teul	&	teul	&	teul\\
4	&	paat	&	faat	&	faat	&	paat	&	faat\\
5	&	liim	&	liim	&	liim	&	liim	&	liim\\
6	&	hohon iid	&	neen	&	neen	&	liim nai niid	&	liim nai ida\\
7	&	hoho ruu	&	hitu	&	hitu	&	liim nai rua	&	liim nai rua\\
8	&	hoho teul	&	ualu	&	ualu	&	liim nai telu	&	liim nai telu\\
9	&	hoho paat	&	sia	&	sian	&	liim nai pata	&	liim nai fata\\
10	&	saguul	&	sakuul	&	sakuul	&	saguul	&	saguul\\
20	&	guul ruu	&	rua nuul	&	rua nulo	&	saguul haet rua	&	haet rua\\
30	&	guul teul	&	teul nuul	&	teul nulo	&	saguul haet teul	&	haet teul\\
100	&	atus iid	&	atus iid	&	atus iid	&	atus iid	&	atus iid\\
		\lspbottomrule
	\end{tabular}
\end{table}

Further diversity within \tsc{Mambae} can be seen by examining
the data in \citet{Fogaca2017}.
Northwest Mambae and Central Mambae retain a genitive suffix
\it{-n} which has been lost in South Mambae; e.g.
Northwest/Central Mambae \it{ate-n}, South Mambae \it{ate} ‘liver’.
Similarly, Northwest Mambae and Central Mambae have a nominal
suffix/enclitic \it{-a} of unclear function; e.g.
Northwest/Central Mambae \it{ai-a}, South Mambae \it{ai} ‘tree’.
\footnote{
		Parallels between Northwest/Central Mambae \it{-a}
		and the \tsc{Rote-Meto} determiner **-a are striking.
		One function of this determiner in \tsc{Rote-Meto} languages
		is to mark definiteness.}

Another point of diversity is seen in the forms of metathesis.
All varieties of \tsc{Mambae} have productive CV → VC metathesis (e.g.
\it{lelo → leol} ‘sun’),
but different processes of  vowel assimilation operate after metathesis.
Thus, for instance, South Mambae from Betano village has lowering
of \it{i → e} after initial \it{a} when metathesis operates,
as seen in *hapuy > *api > \it{aep} ‘fire’.
The varieties of Northwest Mambae from Liquiçá
and Hatulia have taken this further with complete assimilation
of the initial vowel: *api > \it{eep-a} [ɛpa] ‘fire’.
Central Mambae, on the other hand,
does not appear to have any vowel assimilation in this particular
form with all known varieties having \it{aif-a} ‘fire’.

Throughout the remainder of the chapter,
we usually cite examples from only one variety each of Welaun,
Kemak, Tokodede, Northwest Mambae, Central Mambae, and South Mambae.
Unless otherwise stated, Welaun data is from the hamlet of Oele{\Q}u
in the village (\it{desa}) of Bauho /bauhoo/.
Kemak data is from the Kailaku variety,
as spoken in the village of Haikesak/Maumutin.
Tokodede data is from the Liquiçá /likisaa/ variety, which is
the focus of \citet{LSG2006}.
Northwest Mambae data are from the subdistrict of Railaco,
Central Mambae from Liquidoe,
and South Mambae from the village of Betano in the subdistrict of Same.

\scp{Sound changes in all \tsc{Central Timor} languages}{05-02}
The strongest evidence for \tsc{Central Timor} as a distinct
subgroup comes from shared *ŋ > **g.
Proto-Central Timor **g then devoices to \it{k} in Welaun,
Central Mambae everywhere, and word initially
and word finally in some varieties of South Mambae.
Examples of *ŋ > **g are given in \trf{05-04}.

\begin{table}\tcp{\tsc{Central Timor} *ŋ > **g}{05-04}
\begin{adjustbox}{width=\textwidth}
	\begin{threeparttable}\st{2pt}
	\begin{tabular}{lllllll}\lsptoprule
gloss	&	ear	&	name	&	soul	&	ten	&	hear	&	wind\\
PMP	&	*təli\tbr{ŋ}a	&	*\tbr{ŋ}ajan	&	*suma\tbr{ŋ}əd	&	*sa-\tbr{ŋ}a-puluq	&	*də\tbr{ŋ}əR	&	**\tbr{ŋ}əlu\su{ǁ}\\ \midrule
Welaun	&	{\hr}li\tbr{k}a-t	&	{\hr}\tbr{k}alan-aat	&	{\hr}ma\tbr{k}ar-aat	&	{\hr}sa\tbr{k}ulu	&	{\hr}lo\tbr{k}a	&	{\hr\hr}\tbr{k}olu\\
Kemak	&	{\hr}li\tbr{g}a-r	&	{\hr}\tbr{g}ala-r	&	{\hr}sma\tbr{g}ar-aan\su{‡}	&	{\hr}sapulu\su{Δ}	&	{\hr}re\tbr{g}a	&	\hr\cg{(agi)}\\
Tokodede	&	{\hr}teli\tbr{g}i roa\su{†}	&	{\hr}\tbr{g}alaː	&	{\hr}ma\tbr{g}ar	&	{\hr}sa\tbr{g}ulu	&	\cg{(pliʔi)}	&	{\hr\hr}\tbr{g}elu\\
NW. Mambae	&	{\hr}tli\tbr{g}a	&		&	{\hr}sma\tbr{g}a-n	&	{\hr}sa\tbr{g}uul	&	\cg{(pliik)}	&	{\hr\hr}\tbr{g}elo\\
C. Mambae	&	{\hr}ti\tbr{k}a-n	&		&	{\hr}suma\tbr{k}a-n	&	{\hr}sa\tbr{k}uul	&	\cg{(fliik)}	&	{\hr\hr}\tbr{k}iul-a\\
S. Mambae	&	{\hr}teli\tbr{g}a	&	{\hr}\tbr{k}ala	&	{\hr}samaa\tbr{k}	&	{\hr}sa\tbr{g}uul	&	\cg{(fliik)}	&	{\hr\hr}\tbr{k}eul\\
		\lspbottomrule
	\end{tabular}
		\begin{tablenotes}\tnsize
			\item[†]	{Tokodede \it{teligi roa} ‘ear’ is from \citet{Klamer2002},
								the second element \it{roa} means ‘leaf’.}
			\item[‡]	{Kemak \it{smagar-aan} is from Atabae Kemak \citep{Elliot2020}.}
			\item[Δ]	{Kemak \it{sapulu} ‘ten’ is from *sa-ŋa-puluq with loss of
								*ŋa while other languages have loss of *ap.}
			\item[{ǁ}]	{Compare Dhao \it{ŋəlu} ‘wind’ among others (\srf{15-04-02-02}).
								PMP *haŋin > Kemak \it{agi} ‘wind’ has *ŋ > \it{g}.}
		\end{tablenotes}
	\end{threeparttable}
\end{adjustbox}
\end{table}

Within the wider region,
denasalisation of *ŋ is attested in most Austronesian languages
of the Bomberai Peninsula (\srf{11-05-02},
\trf{11-30}) and the \tsc{East Seram} node of \tsc{Seram-Tanimbar-Bomberai}
also have denasalisation of *ŋ in word medial position (\srf{11-02}, \trf{11-08}).

Within the Timor region,
\citet[67--68]{Edwards2021} discusses the change of **ŋ > \it{k}
in \tsc{East Rote} and \tsc{Meto}.
Most instances of \tsc{Rote-Meto} **ŋ affected by this change
are not retentions of PMP *ŋ but were reintroduced after PMP
*ŋ > Proto-Rote-Meto **n.
\citet[67]{Edwards2021} also gives some lexical data which is
exclusively shared between \tsc{Meto}
and \tsc{Central Timor} languages
and attributes this data,
and denasalisation of **ŋ in \tsc{East Rote}
and \tsc{Meto}, to a period of prehistoric contact between \tsc{Meto}
and \tsc{Central Timor} languages.

While there is thus some evidence that denasalisation of **ŋ
in \tsc{Meto} and \tsc{Central Timor} is due to contact between these languages,
it is important to note the change in \tsc{Meto} (and \tsc{East
Rote}) occurred \it{after} the break-up of Proto-Rote-Meto.
Thus, it cannot be identified as a sound change exclusively shared
between these nodes.\footnote{
		The normal reflex of PMP *ŋ in
		Proto-Rote-Meto is **n with 21 examples.
		There are also eight cases of PMP *ŋ = Proto-Rote-Meto **ŋ,
		and four of PMP *ŋ > Proto-Rote-Meto **ŋg.}

In addition to *ŋ > **g,
there are a handful of changes which provide supporting evidence for \tsc{Central Timor}.
These changes are *R > {\0},
**r; *ma-p > **mp; *ə > \it{a} in final syllables;
and perhaps *d > **r.
Each change is discussed in turn.

\tsc{Central Timor} languages show a split of PMP *R.
One set of lexical items has *R > {\0} in all known \tsc{Central
Timor} reflexes, while another set of lexical items attests retention of *R in some languages.
Thus, loss of *R in a subset of lexical items can be identified
as a shared innovation at the level of \tsc{Central Timor}.
Galolen and the \tsc{Eastern Timor} node of \tsc{Timor-Babar}
also have *R > {\0},
though in \tsc{Eastern Timor} languages it is more regular.
Cognate sets which show *R > {\0} in all \tsc{Central Timor}
languages are given in \trf{05-05},
and examples which show retention of *R are given in \trf{05-06}.

\begin{table}[p]\centering\tcp{\tsc{Central Timor} *R > {\0}}{05-05}
%\begin{adjustbox}{width=\textwidth}
	\begin{threeparttable} \st{4.5pt}
	\begin{tabular}{lllllll}\lsptoprule
local gl.	&	{\hr}house	&	{\hr}bone	&	{\hr}bite	&	{\hr}new	&	{\hr}bathe	&	{\hr}day\\
PMP	&	\hr*\tbr{R}umaq	&	\hr*du\tbr{R}i	&	\hr*ka\tbr{R}at	&	\hr*baqə\tbr{R}u	&	\hr*di\tbr{R}us	&	\hr*wa\tbr{R}i\su{‡}\\
P.-C. Timor	&	**uma	&	**dui	&	**kaat	&	**bəu(n)	&	**dius	&	**wairua\\ \midrule
Welaun	&	{\hr\hr}uma	&	{\hr\hr}lui-t	&	{\hr\hr}kaat	&	{\hr\hr}founaan	&	{\hr\hr}diis	&	{\hr\hr}wainrua\\
Kemak	&	{\hr\hr}uma	&	{\hr\hr}rui-r	&	{\hr\hr}aat	&	{\hr\hr}heun	&	{\hr\hr}ruis	&	{\hr\hr}bairua\\
Tokodede	&	{\hr\hr}umo\su{†}	&	{\hr\hr}rui	&	{\hr\hr}kaa	&	{\hr\hr}heuː	&	{\hr\hr}diu	&	{\hr\hr}airuu-n\\
NW. Mambae	&	\hr\cg{(pada)}	&	{\hr\hr}rui-n	&		&	{\hr\hr}heu	&	{\hr\hr}riu	&	{\hr\hr}airua-n\\
C. Mambae	&	\hr\cg{(fada)}	&	{\hr\hr}rui-n	&		&	{\hr\hr}heun	&	{\hr\hr}lariu	&	{\hr\hr}airua-n\\
S. Mambae	&	{\hr\hr}uum	&	{\hr\hr}rui	&	\hr\cg{(hoi)}	&	{\hr\hr}heu	&	{\hr\hr}rio	&	{\hr\hr}airua\\
		\lspbottomrule
	\end{tabular}
		\begin{tablenotes}\tnsize
			\item[†]	{Tokodede \it{umo} is from the Maubara variety.
								Other varieties of Tokodede have \it{soa} ‘house’.}
			\item[‡]	{Reflexes of *waRi are compounds of *waRi + *duha ‘two’.
								Welaun \it{wainrua} and Kemak \it{bairua} are given as ‘day after tomorrow’.
								Tokodede \it{airuu-n}, and \tsc{Mambae} \it{airua-n, airua} are
								given as ‘day before yesterday’.
								Tokodede also has \it{aer-ruu} ‘day after tomorrow’
								and \it{aer} ‘tomorrow’ which may also be cognate with *waRi
								with *R = \it{r}
								and final CV > VC metathesis.
								Another cognate set related to this etymon is *pa-waRi > Welaun
								\it{a-wai}, Kemak \it{pai}, Tokodede \it{pai} ‘dry in sun’.}
		\end{tablenotes}
	\end{threeparttable}
%\end{adjustbox}
\end{table}

\begin{table}[p]\centering\tcp{\tsc{Central Timor} *R = **r}{05-06}
\begin{adjustbox}{width=\textwidth}
	\begin{threeparttable}\st{2pt}
	\begin{tabular}{llllllll}\lsptoprule
local gl.	&	{\hr}root	&	{\hr}brother	&	{\hr}red	&	{\hr}blood	&	{\hr}run	&	{\hr}water	&	{\hr}bamboo\\
PMP	&	\hr*\tbr{R}amut	&	\hr*ña\tbr{R}a	&	\hr*ma-i\tbr{R}aq	&	\hr*da\tbr{R}aq	&	\hr*la\tbr{R}iw	&	\hr*wahiyə\tbr{R}	&	\hr*qau\tbr{R}\\
P.-C. Timor	&	**\tbr{r}amut	&	**na\tbr{r}a	&	**me\tbr{r}a	&	**da\tbr{r}a	&	**pala\tbr{r}i	&	**wee\tbr{r}	&	**au\tbr{r}\\ \midrule
Welaun	&	\hr\cg{(waʔa-n)}	&	\hr\cg{(naʔi)}	&	{\hr\hr}meak	&	{\hr\hr}laa-t	&	{\hr\hr}a-lai	&	{\hr\hr}wee	&	{\hr\hr}au\\
Kemak	&	\hr\cg{(haʔa-n)}	&	{\hr\hr}naa-r	&	{\hr\hr}meak	&	{\hr\hr}raa-r	&	{\hr\hr}pə̆lai	&	{\hr\hr}bia, bea\su{Δ}	&	{\hr\hr}oa\\
Tokodede	&	{\hr\hr}\tbr{r}amut	&	\hr\cg{(ambou)}	&	{\hr\hr}meoː	&	{\hr\hr}raa	&	{\hr\hr}plai	&	{\hr\hr}ee	&	{\hr\hr}betul ae\su{ǁ}\\
NW. Mambae	&	{\hr\hr}\tbr{r}amu-n	&	{\hr\hr}na\tbr{r}a	&	{\hr\hr}me\tbr{r}a	&	{\hr\hr}la\tbr{r}a\su{†}	&		&	{\hr\hr}ee\tbr{r}-a	&	\hr\cg{(betun)}\\
C. Mambae	&	{\hr\hr}\tbr{r}amu-n	&	{\hr\hr}na\tbr{r}a-n	&	{\hr\hr}me\tbr{r}a-n	&	{\hr\hr}la\tbr{l}a	&		&	{\hr\hr}ee\tbr{r}-a	&	\hr\cg{(betun)}\\
S. Mambae	&	\hr\cg{(haa)}	&	{\hr\hr}na\tbr{r}a	&	{\hr\hr}me\tbr{r}a	&	{\hr\hr}la\tbr{r}a	&	{\hr\hr}flae\tbr{r}\su{‡}	&	{\hr\hr}ee\tbr{r}	&	{\hr\hr}oo\tbr{r}\\
		\lspbottomrule
	\end{tabular}
		\begin{tablenotes}\tnsize
			\item[†]	{\tsc{Mambae} has dissimilation of the initial liquid in this
								form; *daRaq > **dara > **rara > \it{lara}.
								Central Mambae from Liquidoe has subsequent assimilation of the
								medial liquid to yield \it{lala}.
								Other varieties of Central Mambae in \citet{Fogaca2017} have
								\it{lara} or \it{lara-n} ‘blood’.}
			\item[‡]	{South Mambae \it{flaer} ‘run’ is a metathesised form.
								Because South Mambae has lowering of /i/ after metathesis when
								the previous vowel is /a/ (e.g.
								*kami \it{ami → aem} ‘we (excl.)’),
								the unmetathesised form of \it{flaer} ‘run’ could be either **flari or **flare.}
			\item[Δ]	{Most varieties of Kemak have \it{bia} ‘water’.
								Marobo and Leosibe Kemak have \it{bia {\tl} bea}.}
			\item[{ǁ}]	{Tokodede \it{betul ae} is glossed as Tetun Dili \it{au
								maus}, English ‘the small bamboo’.
								The term \it{betul} (from PMP *bətuŋ) is given as the generic
								word for bamboo \citep[4]{LSG2006}.}
		\end{tablenotes}
	\end{threeparttable}
\end{adjustbox}
\end{table}

Among the words in \trf{05-06} with retention of *R,
the sound changes are simplest in \tsc{Mambae} and Welaun.
\tsc{Mambae} has *R = \it{r} in all positions,
while Welaun has *R > {\0} in all positions.
The usual pattern in Tokodede is *R > {\0},
with the exception of *Ramut > \it{ramut} ‘root’.
Additionally, *qauR > \it{betul ae} `bamboo' has fronting
of the final vowel along with loss of *R.
This may indicate that *R > {\0} in Tokodede went through an
intermediate glide stage **y.

Kemak has *R > {\0} word medially,
but *R > \it{a} word finally.
Evidence that the final vowel of Kemak \it{bia,
bea} ‘water’ and \it{oa} ‘bamboo’ is indeed a reflex of *R comes
from words with a final double vowel such as *buaq > \it{boo,
buu} ‘betel nut’, *kahu > \it{oo} ‘you (sg.)’,
and *həsi > \it{sii} ‘meat,
flesh’ which do not show final \it{a}.
Another possible example of *R > \it{a} word finally in Kemak
is *niuR > \it{nua,
noa} (Kutubaba Kemak \it{noa} ‘coconut’,
Marobo Kemak \it{noa {\tl} nua} ‘coconut’),
though in this case \tsc{Mambae} has *R > {\0},
as shown by Northwest/Central Mambae \it{noo-a},
South Mambae \it{noo} ‘coconut’.

Another sound change shared between all \tsc{Central Timor} languages
is *ə > \it{a} in most final syllables.
The main exception is when the final consonant is *q,
in which case there is sporadic raising of *ə > \it{e, i}.
Both these sound changes are also found in some \tsc{Timor-Babar}
languages spoken on the Timor mainland.
Examples of *ə in final syllables are given in \trf{05-07}. (See
also \trf{05-10} on \prf{tab:05-10} for another example of *ə > \it{e,
i} /{\gap}q{\#} in reflexes of *daRəq ‘earth’).


\begin{table}\tcp{\tsc{Central Timor} *ə /{\gap}(C){\#}}{05-07}
%\begin{adjustbox}{width=\textwidth}
	\begin{threeparttable}\st{3pt}
	\begin{tabular}{lllllll}\lsptoprule
local gl.	&	teeth	&	black	&	inside	&	fly \cg{(n.)}	&	wash	&	foam\\
PMP	&	*nip\tbr{ə}n	&	*ma-qit\tbr{ə}m	&	*dal\tbr{ə}m\su{†}	&	*lal\tbr{ə}j	&	*bas\tbr{ə}q\su{Δ}	&	*buj\tbr{ə}q\\ \midrule
Welaun	&	\cg{(kiis-aat)}	&	{\hr}met\tbr{a}n	&	{\hr}lam\tbr{a}raan	&	{\hr}lal\tbr{a}n	&	{\hr}sah\tbr{i}	&	{\hr}for\tbr{e}n\\
Kemak	&	{\hr}nip\tbr{a}-r	&	{\hr}met\tbr{a}m	&	{\hr}lar\tbr{a}-n	&	{\hr}lal\tbr{a}	&	{\hr}has\tbr{a}	&	{\hr}bia hul\tbr{a}n\\
Tokodede	&	{\hr}nip\tbr{a}	&	{\hr}met\tbr{a}:	&	{\hr}kai lar\tbr{a}\su{‡}	&	{\hr}lal\tbr{a}	&	{\hr}kahas\tbr{a}	&	{\hr}hur\tbr{i}ː\\
NW. Mambae	&	{\hr}nip\tbr{a}-n	&	{\hr}met\tbr{a}	&	{\hr}ai lal\tbr{a}-n\su{‡}	&		&	{\hr}ha\tbr{a}s	&	\\
C. Mambae	&	{\hr}nif\tbr{a}-n	&	{\hr}met\tbr{a}	&	{\hr}ai lal\tbr{a}-n\su{‡}	&		&	{\hr}ha\tbr{e}s	&	\\
S. Mambae	&	{\hr}nif\tbr{a}	&	{\hr}met\tbr{a}	&	{\hr}lal\tbr{a}	&		&	{\hr}ha\tbr{a}s	&	\\
		\lspbottomrule
	\end{tabular}
		\begin{tablenotes}\tnsize
			\item[†]	{The Kemak,
								Tokodede, and Welaun reflexes all have metathesis of the initial
								and medial consonants: *daləm > **ralam > **laram.
								(Such metathesis may have also taken place in \tsc{Mambae},
								though assimilation of **r > \it{l} means that we cannot be sure.)
								Welaun then has metathesis of the medial
								and final consonant to yield **laram > **lamar.
								This consonant metathesis in Welaun must have been complete before
								the change of non-final **r > \it{l}.
								Final \it{aan} in Welaun is a historic genitive suffix.}
			\item[‡]	{Tokodede \it{kai lara}
								and Northwest/Central Mambae \it{ai lala-n} mean ‘forest’ from ‘tree’ + ‘inside’.
								Kemak also has \it{ai lara-n} ‘forest’.}
			\item[Δ]	{Welaun \it{sahi} ‘wash’ has consonant metathesis.
								Tokodede \it{kahasa} ‘wash’ is from \citet{Klamer2002}.
								The unmetathesised form of Central Mambae \it{haes} ‘wash’ may
								be *[hasi] or *[hase].
								The unmetathesised form of South Mambae \it{haas} ‘wash’ would probably be *[hasa].}
		\end{tablenotes}
	\end{threeparttable}
%\end{adjustbox}
\end{table}

Another change which is probably shared between all \tsc{Central
Timor} languages involves the reflexes of *ma-p.
Welaun has *ma-p > \it{f} while other \tsc{Central Timor} languages
usually have *ma-p > \it{b},
though there are also sporadic cases of *ma-p > \it{p}.
(Instances of *ma-p > \it{p} may simply be reflexes of *p without the prefix *ma-).
Because these reflexes must be kept distinct from the reflexes
of *p (see \trf{05-01}, \prf{tab:05-01}),
they probably point to an intermediate nasal-stop sequence **mp
at the level of Proto-Central Timor.
Examples of PMP *ma-p > **mp in \tsc{Central Timor} are given
in \trf{05-08}.

\begin{table}[p]\centering\tcp{\tsc{Central Timor} *ma-p > **mp}{05-08}
%\begin{adjustbox}{width=\textwidth}
	\begin{threeparttable}\st{2pt}
	\begin{tabular}{llllll}\lsptoprule
local gl.	&	{\hr}hot	&	{\hr}white	&	{\hr}full	&	{\hr}bitter	&	{\hr}sick\\
PMP	&	\hr*\tbr{ma-p}anas	&	\hr*\tbr{ma-p}utiq	&	\hr*\tbr{ma-p}ənuq	&	**(\tbr{ma-})\tbr{p}əju\su{ǁ}	&	**\tbr{ma-p}əjəs\su{¶}\\
P.-C. Timor	&	**\tbr{mp}anas	&	**\tbr{mp}uti	&	**\tbr{mp}ənu	&	**\tbr{mp}əju	&	**\tbr{mp}əjas\\ \midrule
Welaun	&	{\hr\hr}\tbr{f}anas	&	\hr\cg{(baras)}	&	{\hr\hr}\tbr{h}onu\su{Δ}	&	{\hr\hr}\tbr{f}oru	&	{\hr}\tbr{f}oras\\
Kemak	&	{\hr\hr}\tbr{b}ansan\su{‡}	&	{\hr\hr}\tbr{b}uti	&	{\hr\hr}\tbr{b}enu	&	{\hr\hr}\tbr{p}elun	&	\cg{(bə̆ragan)}\\
Tokodede	&	{\hr\hr}\tbr{b}anaː	&	{\hr\hr}\tbr{b}uti	&	{\hr\hr}\tbr{b}enu	&	{\hr\hr}\tbr{p}elu	&	\cg{(groʔa)}\\
NW. Mambae	&	{\hr\hr}\tbr{b}ana	&	{\hr\hr}\tbr{b}uti	&	{\hr\hr}\tbr{b}enu	&		&	\cg{(baan)}\\
C. Mambae\su{†}	&	{\hr\hr}\tbr{p}aan	&	{\hr\hr}\tbr{b}uti	&	{\hr\hr}\tbr{p}eun	&		&	\cg{(paan)}\\
S. Mambae	&	\hr\cg{(brusi)}	&	{\hr\hr}\tbr{b}uti	&	{\hr\hr}\tbr{b}eun	&		&	\cg{(moras)}\\
		\lspbottomrule
	\end{tabular}
		\begin{tablenotes}\tnsize
			\item[†]	{Central Mambae words here are from Laulara.
								The word \it{paan} is given as ‘sick’.}
			\item[‡]	{Kemak \it{bansan} ‘hot’ is from intermediate **banas-aan with
								deletion of the second vowel.
								Marobo Kemak is known to have both \it{bansan} ‘hot’ and \it{banas-aan} ‘heat’.}
			\item[Δ]	{Welaun \it{honu} ‘full’ is from *pənuq without initial *ma-.}
			\item[{ǁ}]	{**ma-pəju ‘bitter’ is from PMP *qapəju ‘gall, gallbladder, bile’.}
			\item[¶]	{**ma-pəjəs is from PMP *hapəjəs ‘smarting, stinging pain’.
								South Mambae \it{moras} ‘sick’ is a borrowing from Tetun,
								as evidenced by irregular *ə > \it{o} rather than expected *ə > \it{e}.}
		\end{tablenotes}
	\end{threeparttable}
%\end{adjustbox}
\end{table}

The reflexes of putative **mp also have to be kept distinct from reflexes of **mb.
While both nasal-stop sequences merge as \it{b} in most \tsc{Central
Timor} languages, in Welaun the reflexes are different with **mb > \it{b} as opposed
to **mp > \it{f.} Examples of **mb are given in \trf{05-09}.

\begin{table}[p]\centering\tcp{\tsc{Central Timor} **mb}{05-09}
%\begin{adjustbox}{width=\textwidth}
	\begin{threeparttable} \st{2pt}
	\begin{tabular}{ll@{}l@{}ll@{}l}\lsptoprule
local gl.	&	{\hr}owner, master	&	{\hr}carrying stick	&	{\hr}betel nut	&	saliva	&	slap\\
PMP	&	\hr*umpu	&	\hr*lemba	&	\hr*buaq	&	**kambeR\su{ǁ}	&	**mbasaR\su{¶}\\
P.-C. Timor	&	**u\tbr{mb}u	&	**kai le\tbr{mb}a\su{‡}	&	**\tbr{mb}ua\su{Δ}	&	**ka\tbr{mb}a	&	**\tbr{mb}asa\\ \midrule
Welaun	&	{\hr\hr}u\tbr{b}un	&	{\hr\hr}ai kale\tbr{b}a	&	{\hr\hr}\tbr{b}ua	&	{\hr\hr}a\tbr{b}a-t	&	{\hr\hr}\tbr{b}asa\\
Kemak	&	{\hr\hr}u\tbr{b}un\su{†}	&	{\hr\hr}ai le\tbr{b}an	&	{\hr\hr}\tbr{b}oo	&	{\hr\hr}a\tbr{b}a-r bea-n\su{†}	&	{\hr\hr}\tbr{b}asa\\
Tokodede	&	{\hr\hr}u\tbr{b}u	&	{\hr\hr}kai le\tbr{b}a	&	{\hr\hr}\tbr{b}uo	&	{\hr\hr}\tbr{b}aa, ka\tbr{b}ee	&	{\hr\hr}k\tbr{b}aas\\
S. Mambae	&	{\hr\hr}u\tbr{b}u	&		&	{\hr\hr}\tbr{b}uu	&	{\hr\hr}a\tbr{b}a	&	{\hr\hr}\tbr{b}aas\\
		\lspbottomrule
	\end{tabular}
		\begin{tablenotes}\tnsize
			\item[†]	{Kemak \it{ubun} ‘owner,
								master’ and \it{aba-r bea-n} ‘saliva’ are from Atabae Kemak \citep{Elliot2020}.}
			\item[‡]	{PCMP *lemba is ‘carry on a shoulder pole’.
								Reflexes here are compounded with reflexes of *kayu ‘wood’
								and refer to the pole with which something is carried over the shoulder.}
			\item[Δ]	{PMP *buaq often develops into a form with an initial nasal-stop
								cluster in languages of Timor when the meaning is ‘betel nut’.
								Compare, for instance, Proto-Rote-Meto *mbuah \citep[270]{Edwards2021}.}
			\item[{ǁ}]	{For **kambeR ‘saliva’ compare Proto-Rote-Meto *kambe,
								Kisar \it{apar,} Kamarian \it{aper} ‘saliva’ among other cognates \citep[192]{Edwards2021}.
								Tokodede \it{baa} is from \citet{LSG2006},
								\it{kabee} is from \citet{Klamer2002}.
								\it{kabee} may a compound of **kaba
								and \it{ee} ‘water’ or a borrowing from Tetun \it{kabeen} ‘saliva’.}
			\item[¶]	{For **mbasaR ‘slap’ compare Proto-Rote-Meto *mbasa,
								Helong \it{papas}, Ili{\Q}uun \it{paas},
								Kisar \it{pahar}, and Sika \it{pahak} among other cognates \citep[263]{Edwards2021}.}
		\end{tablenotes}
	\end{threeparttable}
%\end{adjustbox}
\end{table}

Finally, all \tsc{Central Timor} languages also have *d > **r.
While this results in a merger with those words which retain
*R in \tsc{Mambae}, other languages do not merge *d and *R.
Because of this, we must either assign a value other than [r]
to the Proto-Central Timor reflex of *R (when retained),
and/or we must posit that *d > **r occurred after the break-up
of the group, and would thus provide no evidence for \tsc{Central Timor}.
Whatever hypothesis is taken,
the evidence that *d > **r provides for \tsc{Central Timor} would
be negligible since it is an extremely common change.

Welaun has subsequent *d > **r > \it{l} word initially
and medially, but retains *d > **r = \it{r} word finally,
as shown in \trf{05-10}.
\tsc{Mambae} has occasional instances of *d > **r > \it{l} which
can usually be attributed to assimilation or dissimilation with another adjacent liquid.
Two examples of \tsc{Mambae} *d > \it{l} are *daRaq > **rara
> \it{lara, lala} ‘blood’ (\trf{05-06}, \prf{tab:05-06}),
and *daləm > **rala > \it{lala} ‘inside’ (\trf{05-07}, \prf{tab:05-07}).

\begin{table}\tcp{\tsc{Central Timor} *d}{05-10}
\begin{adjustbox}{width=\textwidth}
	\begin{threeparttable}\st{2pt}
	\begin{tabular}{llllllll}\lsptoprule
local gl.	&	bone	&	earth	&	leaf	&	stand	&	young	&	knee	&	soul\\
PMP	&	*\tbr{d}uRi	&	*\tbr{d}aRəq	&	*\tbr{d}ahun	&	*\tbr{d}iRi	&	*ŋu\tbr{d}a\su{‡}	&	*tuhu\tbr{d}\su{ǁ}	&	*sumaŋə\tbr{d}\\ \midrule
Welaun	&	{\hr}\tbr{l}ui-t	&	{\hr}\tbr{l}ai	&	{\hr}\tbr{l}oon	&	\cg{(a-fii)}	&	{\hr}ku\tbr{l}aʔaan	&	{\hr}tuu\tbr{r}-aat	&	{\hr}maka\tbr{r}-aat\\
Kemak	&	{\hr}\tbr{r}ui-r	&	{\hr}\tbr{r}ae	&	\cg{(ai taha-n)}	&	\cg{(ara)}	&		&	{\hr}tu\tbr{r}u-r	&	{\hr}smaga\tbr{r}-aan\\
Tokodede	&	{\hr}\tbr{r}ui	&	{\hr}\tbr{r}ae	&	{\hr}\tbr{r}oaː	&	{\hr}b\tbr{r}ii\su{Δ}	&		&	\cg{(tunuruː)}	&	{\hr}maga\tbr{r}\\
NW. Mambae	&	{\hr}\tbr{r}ui-n	&	{\hr}\tbr{r}ai-a	&	\cg{(ai nora)}\su{†}	&		&	{\hr}gu\tbr{r}a	&		&	{\hr}smaga-n\\
C. Mambae	&	{\hr}\tbr{r}ui-n	&	{\hr}\tbr{r}ai-a	&	\cg{(nora-n)}\su{†}	&		&	{\hr}ku\tbr{r}a	&		&	{\hr}sumaka-n\\
S. Mambae	&	{\hr}\tbr{r}ui	&	{\hr}\tbr{r}ae	&	\cg{(ai nora)}	&	{\hr}m\tbr{r}ii	&	{\hr}kuu\tbr{r}	&	\cg{(oe tunur)}	&	{\hr}samaak\\
		\lspbottomrule
	\end{tabular}
		\begin{tablenotes}\tnsize
			\item[†]	{Northwest Mambae has \it{ulu laun} ‘head hair’
								and Central Mambae has \it{ulu lahon} ‘head hair’ the second
								elements of which may be irregularly connected with *dahun ‘leaf’ (see \srf{02-03-01-06}).}
			\item[‡]	{The \tsc{Mambae} reflexes mean ‘grass’.
								The reconstructed PMP gloss is ‘young, of plants’.}
			\item[Δ]	{Tokodede \it{brii} ‘stand’ is from \citet{Klamer2002}.}
			\item[{ǁ}]	{Tokodede \it{tunuruː}
								and South Mambae \it{oe tunur} ‘knee’ may be irregularly connected
								with PMP *tuhud, though the source of medial \it{n} is unknown.
								(South Mambae \it{oe tunur} is from \citealt{Klamer2002}).}
		\end{tablenotes}
	\end{threeparttable}
\end{adjustbox}
\end{table}

To summarise, the change of PMP *ŋ > **g can be identified as
a shared innovation which defines the \tsc{Central Timor} subgroup.
While \tsc{Meto} and \tsc{East Rote} have similar **ŋ > \it{k},
this change occurred after the break-up of the \tsc{Rote-Meto}
group and does not affect most instances of PMP *ŋ in \tsc{Rote-Meto}.
Other sound changes shared between \tsc{Central Timor} languages
which provide supporting evidence for this group are the split
of *R > {\0},
**r in the same words,
*ə > \it{a} in word-final syllables,
*ma-p > **mp, and perhaps *d > **r.

\scp{\tsc{Nuclear Central Timor}}{05-03}
Within \tsc{Central Timor} there is evidence for a \tsc{Nuclear
Central Timor} subgroup which contains all languages except Welaun.
Evidence for this subgroup comes from lowering of high vowels
to mid before  final *R,
the merger of **mb/ **mp > \it{b} in most cases (as discussed
above), and penultimate *ə > \it{e}.
Firstly, \tsc{Nuclear Central Timor} languages usually have lowering
of high vowels to mid before final *R,
which was then often lost.
Examples are given in \trf{05-11}.

\begin{table}\tcp{\tsc{Central Timor} *-uR{\#} and *-iR{\#}}{05-11}
\begin{adjustbox}{width=\textwidth}
	\begin{threeparttable}\st{2pt}
	\begin{tabular}{llllll}\lsptoprule
local gl.	&	tail	&	egg	&	grain head	&	dust, ashes\su{Δ}	&	breadfruit\\
PMP	&	*ik\tbr{uR}	&	*qatəl\tbr{uR}	&	*bul\tbr{iR}	&	*dap\tbr{uR}	&	*kul\tbr{uR}\\ \midrule
Welaun	&	{\hr}hiʔ\tbb{u}-n\su{†}	&	{\hr}man tol\tbb{u}-n	&	{\hr}ful\tbb{i}n	&	\cg{(lai lahuk-aan)}	&	{\hr}kul\tbb{u}\\
Kemak	&	{\hr}hiʔ\tbr{o}-r\su{†}	&	{\hr}manu tel\tbr{o}-n	&	{\hr}hul\tbr{e}-n	&	{\hr}rap\tbr{o}, rap\tbr{o} apan	&	{\hr}kul\tbr{o}\\
Tokodede	&	{\hr}ik\tbr{o}	&	{\hr}manu tel\tbr{o}-n	&		&	{\hr}rap\tbr{o}	&	{\hr}kul\tbr{o}r\\
NW. Mambae	&	{\hr}i\tbr{o}-n	&	{\hr}man tel\tbr{o}-n	&		&		&	\\
C. Mambae	&	{\hr}i\tbr{o}-n	&	{\hr}maun tel\tbr{o}-n	&		&		&	\\
S. Mambae\su{†}	&	{\hr}i\tbr{o}	&	{\hr}man tel\tbr{o}	&		&	{\hr}aap rap\tbr{o}	&	\\
		\lspbottomrule
	\end{tabular}
		\begin{tablenotes}\tnsize
			\item[†]	{South Mambae data from Hatu-Udo subdistrict.}
			\item[‡]	{Welaun \it{hiʔu-n}
								and Tokodede \it{hiʔo-r} ‘tail’ have insertion of /h/ /{\#\gap}VʔV.}
			\item[Δ]	{Kemak \it{rapo} is ‘dust; slaked lime,
								calcium’ and \it{rapo apan} is ‘ash in the fireplace’.
								Tokodede \it{rapo} is only given as ‘lime, calcium’ by \citet[23]{LSG2006}.
								South Mambae \it{aap rapo} means ‘ashes’.
								(Initial \it{aap} ‘fire’ is metathesised from **api.) South Mambae
								from Letefoho village has \it{aep ni rapu} with final *u = \it{u} /{\gap}R{\#}.
								The semantic shift from PMP *dapuR ‘hearth’ > ‘calcium’ seems
								motivated by the similarity between the ash in the hearth
								and powdered calcium, as is used for chewing betel nut.}
		\end{tablenotes}
	\end{threeparttable}
\end{adjustbox}
\end{table}

Secondly, all \tsc{Central Timor} languages except Welaun have
penultimate *ə > \it{e}.
Welaun has penultimate *ə > \it{o} in most cases,
which is also the reflex in Tetun (\srf{03-03}).
(One exceptional word in Welaun is *ənəm > \it{inam} ‘six’.)
Examples of *ə are given in \trf{05-12}.

\begin{table}\tcp{\tsc{Central Timor} non-final *ə}{05-12}
\begin{adjustbox}{width=\textwidth}
	\begin{threeparttable}\st{2pt}
	\begin{tabular}{llllllll}\lsptoprule
local gl.	&	three	&	corn\su{†}	&	sun	&	new	&	egg	&	hear	&	flea\\
PMP	&	*t\tbr{ə}lu	&	*z\tbr{ə}lay	&	*qal\tbr{ə}jaw	&	*baq\tbr{ə}Ru	&	*qat\tbr{ə}luR	&	*d\tbr{ə}ŋəR	&	*qatim\tbr{ə}la\\ \midrule
Welaun	&	{\hr}t\tbr{o}lu	&	{\hr}s\tbr{o}le	&	{\hr}l\tbr{o}ro	&	{\hr}f\tbr{o}unaan	&	{\hr}man t\tbr{o}lu-n	&	{\hr}l\tbr{o}ka	&	{\hr}tam\tbr{o}la\\
Kemak	&	{\hr}t\tbr{e}lu	&	{\hr}s\tbr{e}le	&	{\hr}l\tbr{e}lo	&	{\hr}h\tbr{e}un	&	{\hr}manu t\tbr{e}lo-n	&	{\hr}r\tbr{e}ga	&	{\hr}tə̆m\tbr{e}la\\
Tokodede	&	{\hr}t\tbr{e}lu	&	{\hr}s\tbr{e}le	&	{\hr}l\tbr{e}lo	&	{\hr}h\tbr{e}uː	&	{\hr}manu t\tbr{e}lo-n	&	\cg{(pliʔi)}	&	{\hr}tum\tbr{e}l\\
NW. Mambae	&	{\hr}t\tbr{e}ul	&	{\hr}s\tbr{e}la	&	{\hr}l\tbr{e}lon	&	{\hr}h\tbr{e}u	&	{\hr}man t\tbr{e}lo-n	&	\cg{(pliik)}	&	\\
C. Mambae	&	{\hr}t\tbr{e}ul	&	{\hr}s\tbr{e}la	&	{\hr}l\tbr{e}el-a	&	{\hr}h\tbr{e}un	&	{\hr}maun t\tbr{e}lo-n	&	\cg{(fliik)}	&	\\
S. Mambae	&	{\hr}t\tbr{e}ul	&	\cg{(batar)}	&	{\hr}l\tbr{e}ol mata	&	{\hr}h\tbr{e}u	&	{\hr}maun t\tbr{e}lo	&	\cg{(fliik)}	&	\\
		\lspbottomrule
	\end{tabular}
		\begin{tablenotes}\tnsize
			\item[†] {PMP *zəlay is reconstructed with the meaning ‘Job’s tears’.}
		\end{tablenotes}
	\end{threeparttable}
\end{adjustbox}
\end{table}

\tsc{Nuclear Central Timor} languages also share medial *j >
\it{l} in five words, given in \trf{05-13}.
While Welaun also has *j > \it{l} in three of these words,
it has *j > \it{r} adjacent to schwa.
This, combined with the change of Welaun *d > **r > \it{l},
indicates that Welaun *j > \it{l} in these words may have gone
though intermediate **d and/or **r,
in which case the change of *j > \it{l} would not be shared between
all \tsc{Central Timor} languages.\footnote{
		Note, however, that Welaun has *d > \it{l} adjacent to schwa as shown
		by *dəŋəR > \it{loka} ‘hear’.
		The only other known cases of *d adjacent to schwa are *dəpa
		> \it{roʔa} ‘stretch out the arm,
		arm-span’, which may be borrowed from West Tetun \it{roʔa},
		and *sumaŋəd > \it{makar-aat} ‘soul’,
		for which *d is word final.}

\begin{table}\tcp{\tsc{Nuclear Central Timor} *j > \it{l}}{05-13}
%\begin{adjustbox}{width=\textwidth}
	\begin{threeparttable} %\st{2pt}
	\begin{tabular}{llllll}\lsptoprule
local gl.	&	ySib\su{†}	&	nose	&	how much	&	name	&	bitter\\
PMP	&	*hua\tbr{j}i	&	*i\tbr{j}uŋ	&	*pi\tbr{j}a	&	*ŋa\tbr{j}an	&	**ma-pə\tbr{j}u\su{Δ}\\	\midrule
Welaun	&	{\hr}wa\tbr{l}i-n	&	{\hr}i\tbr{l}uk-aat	&	{\hr}hi\tbr{r}a	&	{\hr}ka\tbr{l}an-aat	&	{\hr\hr}fo\tbr{r}u\\
Kemak	&	\hp{*w}a\tbr{l}i-r	&	{\hr}i\tbr{l}u-r	&	{\hr}pi\tbr{l}a	&	{\hr}ga\tbr{l}a-r	&	{\hr\hr}pe\tbr{l}un\\
Tokodede	&	\hp{*w}a\tbr{l}i	&	{\hr}i\tbr{l}u	&	{\hr}pii\tbr{l}	&	{\hr}ga\tbr{l}aː	&	{\hr\hr}pe\tbr{l}u\\
NW. Mambae	&	\hp{*w}a\tbr{l}i-n	&	{\hr}i\tbr{l}u-n	&	{\hr}arpii\tbr{l}\su{‡}	&		&	\\
C. Mambae	&	\hp{*w}a\tbr{l}i-n	&	{\hr}i\tbr{l}u-n	&	{\hr}arfii\tbr{l}\su{‡}	&		&	\\
S. Mambae	&	\hp{w}\cg{(kau)}	&	{\hr}i\tbr{l}u	&	{\hr}arpii\tbr{l}\su{‡}	&	{\hr}ka\tbr{l}a	&	\\
		\lspbottomrule
	\end{tabular}
		\begin{tablenotes}\tnsize
			\item[†]	{Reflexes of *huaji mean ‘same sex younger sibling’ in \tsc{Central Timor} languages.}
			\item[‡]	{The Mambae words \it{arpiil, arfiil} mean ‘when?’.}
			\item[Δ]	{**ma-pəju ‘bitter’ is from PMP *qapəju ‘gall, gallbladder, bile’.}
		\end{tablenotes}
	\end{threeparttable}
%\end{adjustbox}
\end{table}

Apart from the five words which have *j > \it{l} in \tsc{Nuclear
Central Timor}, the other known reflexes of *j are: PMP *ma\tbr{j}a ‘fry’ > Welaun
\it{ma\tbr{r}a}; Kemak \it{ma\tbr{l}an}; South Mambae **masa
> \it{maa\tbr{s}}, *bu\tbr{j}əq ‘foam’ > Welaun \it{fo\tbr{r}en}; Kemak \it{bia hu\tbr{l}an}; Tokodede \it{hu\tbr{r}iː},
*pusə\tbr{j} ‘navel; centre’ > Welaun \it{husa\tbr{r}-aat};
Atabae Kemak \it{pus\tbr{r}-aar}; Tokodede \it{kupusa};
South Mambae \it{fusu}, and probably *qulə\tbr{j} > Kemak \it{ula\tbr{l}}, Tokodede
\it{ula} ‘snake’, though this last word may be from *hulaR.

\scp{Other sound correspondences}{05-04}
PMP *z has a number of reflexes in the \tsc{Central Timor} languages
with *z > \it{s} being the most common reflex.
Welaun has *z > \it{s,
d}; Kemak has *z > \it{s}; Tokodede has *z > \it{s,
r}; and \tsc{Mambae} has *z > \it{s} with one possible instance
of *z > \it{r}.
Examples are given in  \trf{05-14}.

\begin{table}\tcp{\tsc{Central Timor} *z}{05-14}
%\begin{adjustbox}{width=\textwidth}
	\begin{threeparttable}\st{2.25pt}
	\begin{tabular}{llllllll}\lsptoprule
local gl.	&	path	&	far	&	corn	&	rain	&	ladder	&	bad	&	point\\
PMP	&	*\tbr{z}alan	&	*\tbr{z}auq	&	*\tbr{z}əlay	&	*qu\tbr{z}an	&	*haRə\tbr{z}an	&	*\tbr{z}aqat	&	*tu\tbr{z}uq\\ \midrule
Welaun	&	{\hr}\tbr{s}alan	&	{\hr}lai\tbr{s}oo\su{†}	&	{\hr}\tbr{s}ole	&	{\hq}u\tbr{s}an	&	{\hr}o\tbr{s}an	&	{\hr}\tbr{d}aat	&	{\hr}a-tu\tbr{d}u\\
Kemak	&	{\hr}\tbr{s}ala	&	\cg{(sə̆goon)}	&	{\hr}\tbr{s}ele	&	{\hq}u\tbr{s}a	&		&	\cg{(mə̆ragan)}	&	\cg{(suku)}\\
Tokodede	&	\cg{(daa)}	&	{\hr}k\tbr{r}oo	&	{\hr}\tbr{s}ele	&	{\hq}u\tbr{r}a	&		&	\cg{(klao)}	&	{\hr}t(u)nu\tbr{s}u\\
NW. Mambae	&	\cg{(daa)}	&	\cg{(room)}\su{‡}	&	{\hr}\tbr{s}ela	&	{\hq}u\tbr{s}a	&		&	\cg{(klao)}	&	\\
C. Mambae	&	\cg{(taan)}	&	{\hr}\tbr{r}oo	&	{\hr}\tbr{s}ela	&	{\hq}u\tbr{s}a	&		&	\cg{(soko)}	&	\\
S. Mambae	&	{\hr}\tbr{s}aal	&	{\hr}\tbr{r}oo	&	\cg{(batar)}	&	{\hq}uu\tbr{s}	&		&	\cg{(klao)}	&	\\
		\lspbottomrule
	\end{tabular}
		\begin{tablenotes}\tnsize
			\item[†]	{Initial \it{lai} in Welaun \it{laisoo} ‘far’ may be from *daRəq > \it{lai} ‘earth’.}
			\item[‡]	{Final \it{m} in Northwest Mambae \it{room} indicates that
								Central Mambae and South Mambae \it{roo} may not be a reflex of *zauq.
								Compare also Northwest Mambae from Hatulia \it{roam} ‘far’.}
		\end{tablenotes}
	\end{threeparttable}
%\end{adjustbox}
\end{table}

PMP *b undergoes lenition in all \tsc{Central Timor} languages.
Kemak, \tsc{Mambae}, Liquiçá Tokodede
and Bazartete Tokodede have *b > \it{h} in all positions.
Some varieties of Kemak then have sporadic initial *h > {\0}.
Maubara Tokodede has *b > \it{v} in all positions.
Welaun has initial *b > \it{f}
and medial *b > \it{h}.
Examples are given in \trf{05-15}.
Two examples from Welaun which are not identical to Tetun
and establish the regularity of the pattern are *bəqbəq ‘mouth’
> \it{foha-t} ‘voice’ and *bubuŋ > \it{fuhuk} ‘ridgepole’.

\begin{table}[p]\centering\tcp{\tsc{Central Timor} *b}{05-15}
%\begin{adjustbox}{width=\textwidth}
	\st{3.5pt}
	\begin{tabular}{lllllll}\lsptoprule
local gl.	&	pig	&	stone	&	star	&	fruit	&	moon	&	body hair\\
PMP	&	*\tbr{b}a\tbr{b}uy	&	*\tbr{b}atu	&	*\tbr{b}ituqən	&	*\tbr{b}uaq	&	*\tbr{b}ulan	&	*\tbr{b}ulu\\ \midrule
Welaun	&	{\hr}\tbr{f}a\tbr{h}i	&	{\hr}\tbr{f}atu	&	{\hr}\tbr{f}itun	&	{\hr}ai \tbr{f}uan	&	{\hr}\tbr{f}ulan	&	{\hr}\tbr{f}ulu-t\\
Kutubaba Kemak	&		&	{\hr}\tbr{h}atu	&	{\hr}\tbr{h}itu	&	{\hr}\tbr{h}uaŋ	&	{\hr}\tbr{h}ula	&	{\hr}\tbr{h}ulu-ŋ\\
Kailaku Kemak	&	{\hh}a\tbr{h}i	&	{\hh}atu	&	{\hh}itu	&	{\hr}ai \tbr{h}uan	&	{\hr}\tbr{h}ula	&	{\hr}\tbr{h}ulu-r\\
Liquiçá Tokodede	&	{\hr}\tbr{h}a\tbr{h}i	&	{\hr}\tbr{h}atu	&	{\hr}\tbr{h}itu	&	{\hr}\tbr{h}uu	&	{\hr}\tbr{h}ula	&	{\hr}\tbr{h}uluː\\
Maubara Tokodede	&	{\hr}\tbr{v}a\tbr{v}i	&	{\hr}\tbr{v}atu	&	{\hr}\tbr{v}itu	&	{\hr}\tbr{v}uu	&	{\hr}\tbr{v}ula	&	{\hr}\tbr{v}uluː\\
NW. Mambae	&	{\hr}\tbr{h}ee\tbr{h}-a	&	{\hr}\tbr{h}ato	&	{\hr}\tbr{h}iut-a	&	{\hr}ai \tbr{h}uan	&	{\hr}\tbr{h}ula	&	{\hr}\tbr{h}ulu-n\\
C. Mambae	&	{\hr}\tbr{h}ae\tbr{h}-a	&	{\hr}\tbr{h}aut-a	&	{\hr}\tbr{h}iut-a	&	{\hr}ai \tbr{h}uan	&	{\hr}\tbr{h}ula	&	{\hr}\tbr{h}ulu-n\\
S. Mambae	&	{\hr}\tbr{h}ae\tbr{h}	&	{\hr}\tbr{h}aut	&	{\hr}\tbr{h}iut	&	{\hr}ai \tbr{h}ua	&	{\hr}\tbr{h}ulai	&	{\hr}\tbr{h}ulu\\
		\lspbottomrule
	\end{tabular}
%\end{adjustbox}
\end{table}

Word initially PMP *k > {\0} in all \tsc{Central Timor} languages
except Tokodede where it is retained as \it{k}.
Word medially *k > \it{ʔ} in Welaun; *k > \it{ʔ,} {\0} in Kemak;
*k > \it{ʔ, k} in Tokodede; and *k > {\0} in most \tsc{Mambae} words,
though South Mambae retains *k = \it{k} in at least one word.
Examples of PMP *k are given in \trf{05-16}.

\begin{table}[p]\centering\tcp{\tsc{Central Timor} *k}{05-16}
\begin{adjustbox}{width=\textwidth}
	\begin{threeparttable}\st{2pt}
	\begin{tabular}{llllllll}\lsptoprule
local gl.	&	wood, tree	&	\tsc{2sg}	&	head louse	&	tail	&	steal	&	ascend	&	elbow\\
PMP	&	*\tbr{k}ahiw	&	*\tbr{k}ahu	&	*\tbr{k}utu	&	*i\tbr{k}uR	&	*na\tbr{k}aw	&	*sa\tbr{k}ay	&	*si\tbr{k}u\\ \midrule
Welaun	&	{\hr}ai	&	{\hr}oo	&	{\hr}utu	&	{\hr}hi\tbr{ʔ}u-n\su{†}	&	{\hr}a-na\tbr{ʔ}o	&	{\hr}sa\tbr{ʔ}e	&	{\hr}si\tbr{ʔ}u-t\\
Kemak	&	{\hr}ai	&	{\hr}oo	&	{\hr}utu	&	{\hr}hi\tbr{ʔ}o-r\su{†}	&	{\hr}pə̆nao	&	{\hr}sae	&	{\hr}si\tbr{ʔ}u-r\\
Tokodede	&	{\hr}\tbr{k}ai	&	{\hr}\tbr{k}oo	&	{\hr}\tbr{k}utu	&	{\hr}i\tbr{k}o	&	{\hr}mana\tbr{ʔ}o	&	{\hr}sa\tbr{ʔ}e	&	\\
NW. Mambae	&	{\hr}ai-a	&		&	{\hr}uut-a	&	{\hr}io-n	&	{\hr}pnao	&	{\hr}leol sae-n\su{‡}	&	{\hr}siu-n\\
C. Mambae	&	{\hr}ai-a	&	{\hr}oo	&	{\hr}utu-n	&	{\hr}io-n	&	{\hr}fnao	&	{\hr}leel-a sae\su{‡}	&	{\hr}lima siu-n\\
S. Mambae	&	{\hr}ai	&	{\hr}oo	&	{\hr}uut	&	{\hr}io	&	{\hr}fnao-tee	&	{\hr}leol sae\su{‡}	&	{\hr}lima si\tbr{k}u\\
		\lspbottomrule
	\end{tabular}
		\begin{tablenotes}\tnsize
			\item[†]	{Welaun \it{hiʔu-n}
								and Tokodede \it{hiʔo-r} ‘tail’ have insertion of /h/ /{\#\gap}VʔV.}
			\item[‡]	{\tsc{Mambae} forms given here mean ‘east’, literally ‘sun rise’.}
		\end{tablenotes}
	\end{threeparttable}
\end{adjustbox}
\end{table}

The usual reflexes of word-final vowel-glide sequences are *-ay
> \it{e} and *-aw > \it{o},
though there are also sporadic reflexes of *-aw/-*ay > \it{a} in \tsc{Mambae}.
Additionally, all \tsc{Nuclear Central Timor} languages have
*-aw > \it{a} in reflexes of *lakaw ‘go’,
with Welaun showing *-aw > \it{a {\tl} o} in this word.
Examples are given in \trf{05-17}.

\begin{table}[p]\centering\tcp{\tsc{Central Timor} *-ay and *-aw}{05-17}
\begin{adjustbox}{width=\textwidth}
	\st{2.5pt}
	\begin{tabular}{lllll|llll}\lsptoprule
local gl.	&	liver	&	die	&	snake	&	corn	&	steal	&	sun	&	rat	&	go\\
PMP	&	*qat\tbr{ay}	&	*mat\tbr{ay}	&	*nip\tbr{ay}	&	*zəl\tbr{ay}	&	*nak\tbr{aw}	&	*qaləj\tbr{aw}	&	*balab\tbr{aw}	&	*lak\tbr{aw}\\ \midrule
Welaun	&	{\hr}at\tbr{e}-t	&	{\hr}mat\tbr{e}	&	\cg{(saba)}	&	{\hr}sol\tbr{e}	&	{\hr}a-naʔ\tbr{o}	&	{\hr}lor\tbr{o}	&	{\hr}lah\tbr{o}	&	{\hr}laʔ\tbr{a}, laʔ\tbr{o}\\
Kemak	&	{\hr}at\tbr{e}-r	&	{\hr}mat\tbr{e}	&	{\hr}nip\tbr{e}	&	{\hr}sel\tbr{e}	&	{\hr}pə̆na\tbr{o}	&	{\hr}lel\tbr{o}	&	{\hr}lah\tbr{o}	&	{\hr}la\tbr{a}\\
Tokodede	&	{\hr}at\tbr{e}	&	{\hr}mat\tbr{e}	&	{\hr}nip\tbr{e}	&	{\hr}sel\tbr{e}	&	{\hr}manaʔ\tbr{o}	&	{\hr}lel\tbr{o}	&		&	{\hr}la\tbr{a}\\
NW. Mam.	&	{\hr}at\tbr{e}-n	&	{\hr}me\tbr{e}t	&	\cg{(sokai-a)}	&	{\hr}sel\tbr{a}	&	{\hr}pna\tbr{o}	&	{\hr}lel\tbr{o}n	&	{\hr}lah\tbr{o}	&	\\
C. Mambae	&	{\hr}at\tbr{e}-n	&	{\hr}ma\tbr{e}t	&	\cg{(solai-a)}	&	{\hr}sel\tbr{a}	&	{\hr}fna\tbr{o}	&	{\hr}leel-a	&	{\hr}lah\tbr{a}	&	\\
S. Mambae	&	{\hr}at\tbr{e}	&	{\hr}ma\tbr{e}t	&	\cg{(sohai)}	&	\cg{(batar)}	&	{\hr}fna\tbr{o}-tee	&	{\hr}le\tbr{o}l mata	&	{\hr}la\tbr{o}h	&	{\hr}la\tbr{a}\\
		\lspbottomrule
	\end{tabular}
\end{adjustbox}
\end{table}

PMP *w is usually retained unchanged word initially in Welaun
with one example of *w > \it{b} medially in *sawa > \it{saba} ‘snake’.
PMP *w is also usually lost in Tokodede and \tsc{Mambae}.
Kemak usually has *w > \it{b},
though *wakaR ‘roots’ > \it{haʔa-n} has *w > {\0} with subsequent
epenthesis of \it{h} /{\#}{\gap}VʔV (compare *ikuR > \it{hiʔo-r} ‘tail’).
Examples are given in \trf{05-18}.

\begin{table}\tcp{\tsc{Central Timor} *w}{05-18}
%\begin{adjustbox}{width=\textwidth}
	\begin{threeparttable} \st{3.5pt}
	\begin{tabular}{llllll}\lsptoprule
local gl.	&	water	&	roots	&	eight	&	day	&	nine\\
PMP	&	*\tbr{w}ahiyəR	&	*\tbr{w}akaR	&	*\tbr{w}alu	&	*\tbr{w}aRi\su{Δ}	&	*si\tbr{w}a\\ \midrule
Welaun	&	{\hr}\tbr{w}ee	&	{\hr}\tbr{w}aʔa-n	&	{\hr}\tbr{w}alu	&	{\hr}\tbr{w}ainrua	&	{\hr}si\tbr{w}i\\
Kemak	&	{\hr}\tbr{b}ia, \tbr{b}ea\su{†}	&	{\hr}haʔa-n\su{‡}	&	{\hr}\tbr{b}alu	&	{\hr}\tbr{b}airua	&	{\hr}si\tbr{b}e\\
Tokodede	&	{\hr}ee	&	\cg{(ramut)}	&	\cg{(hoho telu)}	&	{\hr}airuu-n	&	\cg{(hoho paat)}\\
NW. Mambae	&	{\hr}eer-a	&	\cg{(ramu-n)}	&	\cg{(hoho teul)}	&	{\hr}airua-n	&	\cg{(hoho paat)}\\
C. Mambae	&	{\hr}eer-a	&	\cg{(ramu-n)}	&	{\hr}\tbr{u}alu\su{ǁ}	&	{\hr}airua-n	&	{\hr}sia\su{ǁ}\\
S. Mambae	&	{\hr}eer	&	{\hr}haa	&	\cg{(liim nai telu)}	&	{\hr}airua	&	\cg{(liim nai fata)}\\
		\lspbottomrule
	\end{tabular}
		\begin{tablenotes}\tnsize
			\item[†]	{Most varieties of Kemak have \it{bia} ‘water’.
								Marobo and Leosibe Kemak have \it{bia {\tl} bea}.}
			\item[‡]	{Kemak \it{haʔa-n} ‘roots’ has insertion of /h/ /{\#\gap}VʔV.
								South Mambae \it{haa} ‘roots’ probably also shows this process
								with subsequent medial **ʔ > {\0}.}
			\item[Δ] 	{See note to \trf{05-05} (\prf{tab:05-05}) for the glosses
								of reflexes of *waRi given here.}
			\item[{ǁ}]	{Central Mambae \it{ualu} ‘eight’
								and \it{sia} ‘nine’ may be borrowings from eastern Tetun.}
		\end{tablenotes}
	\end{threeparttable}
%\end{adjustbox}
\end{table}

\largerpage
Some varieties of Kemak,
including Saneri Kemak and Marobo Kemak have *t > \it{ʧ} before high vowels.
Saneri Kemak also appears to have *t > \it{ʧ} word finally after
high vowels, though there is only one example in our data.
Examples are given in \trf{05-19}.

\begin{table}
	\centering\tcp{Kemak *t > \it{ʧ} /V\tsc{[+high]}}{05-19}
	\st{2.25pt}\begin{tabular}{llllllll}\lsptoprule
local gl.			&	head louse				&	roast						&	skin									&	white						&	stone						&	old								&	grass\\
PMP						&	*ku\tbr{t}u				&	*\tbr{t}unu			&	*kuli\tbr{t}					&	*ma-pu\tbr{t}iq	&	*ba\tbr{t}u			&										&	\\ \midrule
Kailaku Kemak	&	\hp{*k}u\tbr{t}u	&									&	\hp{*k}uli\tbr{t}-iir	&	{\hr}bu\tbr{t}i	&	\hp{*b}a\tbr{t}u&	{\hr}\tbr{t}uman	&	{\hr}suu\tbr{t}\\
Marobo Kemak	&	\hp{*k}u\tbr{ʧ}u	&	{\hr}\tbr{ʧ}unu	&	\hp{*k}uli\tbr{ʧ}-iir	&	{\hr}bu\tbr{ʧ}i	&	\hp{*b}a\tbr{ʧ}u&	{\hr}\tbr{ʧ}uman	&	{\hr}suu\tbr{t}\\
Saneri Kemak	&	\hp{*k}u\tbr{ʧ}u-r&	{\hr}\tbr{ʧ}unu	&	\hp{*k}ul\tbr{ʧ}i-r		&	{\hr}bu\tbr{ʧ}i	&	\hp{*b}a\tbr{ʧ}u&	{\hr}\tbr{ʧ}uman	&	{\hr}suu\tbr{ʧ}\\
		\lspbottomrule
	\end{tabular}
\end{table}

Finally, Tokodede and some varieties of Kemak have loss of final
*n with compensatory lengthening of the previous vowel.
Some other varieties of Kemak have *n > \it{ŋ} syllable finally,
and this may represent an intermediate stage before loss of the nasal.
Thus, the full pathway may have been *Vn > **Vŋ > **Vː.
This process is attested in Tokodede, Marobo Kemak,
Lemia Kemak, Diirbatin Kemak, and Atsabe Kemak.
Examples are given in \trf{05-20}.

\begin{table}
	\centering\tcp{Kemak and Tokodede *n > {\0} and compensatory lengthening}{05-20}
	\begin{adjustbox}{width=\textwidth}
	\st{2pt}\begin{tabular}{lllllll}\lsptoprule
local gl.	&	new	&	hot	&	body hair	&	bark	&	straight	&	big\\
PMP	&	*baqəRu	&	*ma-panas	&	*bulu	&	*kayu + *kulit	&		&	\\ \midrule
Leosibe Kemak	&	{\hr}heu\tbr{n}	&	{\hr}bansa\tbr{n}	&	{\hr}hulu-r	&	{\hr}ai ulti-\tbr{n}	&	{\hr}loso\tbr{n}	&	{\hr}bote\tbr{n}\\
Leolima Kemak	&	{\hr}heu\tbr{ŋ}	&	{\hr}baŋsa\tbr{ŋ}	&	{\hr}hulu-\tbr{ŋ}	&	{\hr}au ulti\tbr{ŋ}	&	{\hr}loso\tbr{ŋ}	&	{\hr}bote\tbr{ŋ}\\
Lemia Kemak	&	{\hr}heu\tbr{ː}	&	{\hr}bansa\tbr{ː}	&	{\hr}hulu\tbr{ː}	&	{\hr}ai ulit	&	{\hr}loso\tbr{ː}	&	{\hr}bote\tbr{ː}\\
Diirbatin Kemak	&	{\hr}heu\tbr{ː}	&	{\hr}bansa\tbr{ː}	&	{\hr}lima-r hulu\tbr{ː}	&	{\hr}ai ulti\tbr{ː}	&	{\hr}loso\tbr{ː}	&	{\hr}bote\tbr{ː}\\
Tokodede	&	{\hr}heu\tbr{ː}	&	{\hr}bana\tbr{ː}	&	{\hr}hulu\tbr{ː}	&	{\hr}kuluta\tbr{ː}	&		&	\cg{(aʔe)}\\
		\lspbottomrule
	\end{tabular}
	\end{adjustbox}
\end{table}


\newpage
\scp{\tsc{Central Timor} classification}{05-05}
\vspace{5mm}

\dirtree{%This  is needed
.1 {\tsc{Central Timor} *ŋ > **g; *R > *R, {\0}; *ə > \it{a} /σ\und{σ}{\#}; *ma-p > **mp; *d > **r; \newline*q/h > {\0} (*p = **p)}.
.2 {Welaun [wlh] **g > \it{k}; *ə > \it{o}; **r > \it{l}; **mp/*b- > \it{f} /{\#}{\gap};
								 *z > \it{s, d}; *p-> \it{h} /{\#\gap}; *-b- > \it{h} /V{\gap}V;
								 *j > \it{r, l}; *R > {\0}; **mb > \it{b}}.
.2 {\tsc{Nuclear Central Timor} *V\tsc{[+high]} > V\tsc{[+mid]} /{\gap}*R{\#}; *ə > \it{e}; \newline\hp{**}**mp/**mb > \it{b}}.
		.3 {Kemak [kem] *w > \it{b}; *z > \it{s}; *b > \it{h}; *R > \it{a} /V{\gap\#}; *R > {\0} V{\gap}V}.
		.3 {Tokodede [tkd] *w > {\0}; *z > \it{s, r;} **r/*R > \it{r} (*k = \it{k} /{\#\gap})}.
				.4 {Liquiçá Tokodede *b > \it{h}}.
				.4 {Maubara Tokodede *b > \it{v}}.
		.3 {\tsc{Mambae} [mgm] *w > {\0}; *z > \it{s}; **r/*R > \it{r}; *b > \it{h}}.
				.4 {Northwest Mambae [{-}{-}{-}]}.
				.4 {Central Mambae [{-}{-}{-}] **g > \it{k}; *p > \it{f}}.
				.4 {South Mambae [{-}{-}{-}] **-g > -\it{k} /{\gap\#}; *p > \it{p {\tl} f} (variety dependent)}.
}%Ensure there is no carriage return before this }
